\chapter{Profil organizacyjny, gospodarczy i otoczenie fizyczne}\label{ch:organizacja}

% ==============================================================================
% ROZDZIAŁ 2 - WSKAZÓWKI I TO-DO DLA AGENTA PISZĄCEGO:
% Rola: Osoba 1 (Warstwa organizacyjna, gospodarcza i fizyczna L1-L2).
% Cel: Zarysowanie struktury D-INFK, budynków (CAB), grantów SNSF i przetargów IT.
% Anonimizacja: Tylko oficjalne kierownictwo jawne; reszta personelu przez \anon{...}.
% Objętość docelowa: ~7-8 stron.
% ==============================================================================

\section{Struktura organizacyjna i zarządzanie wydziałem D-INFK}\label{sec:struktura-dinfk}

Analiza struktury organizacyjnej instytucji akademickiej stanowi fundamentalny etap rekonesansu w ramach Filaru I przyjętej metodyki OSINT. Pozwala ona na zidentyfikowanie łańcucha decyzyjnego, podziału kompetencji, nieformalnych zależności oraz kluczowych punktów zaufania wewnątrz organizacji. W przypadku Wydziału Informatyki Politechniki Federalnej w Zurychu (\textit{Departement Informatik - D-INFK}) dekompozycja struktury zarządczej umożliwia nie tylko zrozumienie obiegu informacji formalnych, lecz przede wszystkim dostarcza danych niezbędnych do wiarygodnego modelowania wektorów socjotechnicznych opartych na autorytecie instytucjonalnym.

\subsection{Władze wydziału i organy decyzyjne}\label{subsec:wladze-wydzialu}

W strukturze akademickiej i administracyjnej Politechniki Federalnej w Zurychu Wydział Informatyki (\textit{D-INFK}) funkcjonuje jako jednostka o wysokim stopniu autonomii wykonawczej, podlegając bezpośrednio nadrzędnej Radzie Uczelni (\textit{Schulleitung}) oraz federalnej Radzie Politechnik (\textit{ETH-Rat}). Organem kierowniczym odpowiedzialnym za bieżące zarządzanie strategiczne, operacyjne i naukowe wydziału jest Rada Wydziału (\textit{Departementsleitung - DL}) \cite{eth_dinfk_org}. Działalność kierownictwa jest wspierana przez dwa kluczowe organy kolegialne:
\begin{itemize}
    \item \textbf{Konferencję Wydziału (\textit{Departementskonferenz - DK}):} Najwyższe ciało plenarne zrzeszające wszystkich profesorów zwyczajnych, przedstawicieli kadry naukowej, personelu techniczno-administracyjnego oraz studentów. Konferencja odpowiada za uchwalanie programów studiów, regulaminów wewnętrznych oraz opiniowanie wniosków nominacyjnych.
    \item \textbf{Komitet Wydziałowy (\textit{Departementsausschuss - DA}):} Węższy komitet doradczo-wykonawczy, koordynujący realizację uchwał DK oraz opiniujący kluczowe decyzje finansowo-lokalowe przed ich wdrożeniem.
\end{itemize}

W ramach pasywnego rekonesansu OSINT przeprowadzono ekstrakcję danych ze struktury oficjalnego portalu wydziału (\file{inf.ethz.ch}), co pozwoliło na precyzyjną identyfikację imiennego składu Rady Wydziału (DL), zaprezentowanego na rysunku \ref{fig:dinfk-board} oraz w tabeli \ref{tab:wladze-dinfk}.

\Figure[0.65]{dinfk_board.jpg}{Skład zarządu wydziału D-INFK (\textit{Departementsleitung}) opublikowany na oficjalnym portalu uczelni}{fig:dinfk-board}

\begin{table}[H]
    \centering
    \caption{Skład Rady Wydziału Informatyki ETH Zürich (\textit{Departementsleitung - DL})}\label{tab:wladze-dinfk}
    \tablefont
    \begin{tabularx}{\linewidth}{L{0.28\linewidth} L{0.32\linewidth} Y}
        \toprule
        \hcell{Imię i nazwisko} & \hcell{Funkcja formalna (DL)} & \hcell{Rola operacyjna i znaczenie instytucjonalne} \\
        \midrule
        Prof. Thomas Hofmann & Head of Department & Strategiczne zarządzanie wydziałem, nadzór nad budżetem i reprezentacja zewnętrzna. \\
        Prof. Srdjan Čapkun & Deputy Head of Department & Zastępca szefa wydziału, koordynacja inicjatyw badawczych w obszarze bezpieczeństwa. \\
        Prof. Adrian Perrig & Director of Studies & Nadzór nad tokiem studiów, sprawami studenckimi, egzaminami i kadrą dydaktyczną. \\
        Prof. Valentina Boeva & Deputy Director of Studies & Zastępca dyrektora ds. dydaktyki, koordynacja międzynarodowych programów studiów. \\
        Prof. Stelian Coros & Członek zarządu (Board Member) & Reprezentacja instytutów interaktywnych systemów inteligentnych. \\
        Prof. Siyu Tang & Członek zarządu (Board Member) & Reprezentacja instytutów obliczeń wizualnych i grafiki komputerowej. \\
        Barbara Pfändner & Managing Director & Dyrektor zarządzająca; operacyjna administracja kadr, finansów i infrastruktury. \\
        Stefan Walter & Head IT Services (ISGINF) & Szef wydziałowej grupy wsparcia teleinformatycznego (ISG D-INFK). \\
        \bottomrule
    \end{tabularx}
\end{table}

\lead{Znaczenie wywiadowcze i modelowanie wektorów socjotechnicznych (Actionable Intelligence)}

Identyfikacja ścisłego kierownictwa wydziału dostarcza kluczowych przesłanek dla konstrukcji realistycznych scenariuszy socjotechnicznych, opisanych szczegółowo w rozdziale \ref{ch:socjotechnika}. W inżynierii społecznej mechanizm uległości wobec autorytetu (\textit{Authority Bias}) stanowi jeden z najskuteczniejszych wektorów manipulacji w środowiskach akademickich i korporacyjnych:
\begin{enumerate}
    \item \textbf{Pretexting dziekański (Prof. Thomas Hofmann):} Wykorzystanie tożsamości szefa wydziału umożliwia formułowanie komunikatów o charakterze bezwzględnego priorytetu organizacyjnego (np. rzekome pilne ankiety akredytacyjne, zarządzenia kryzysowe czy wymogi sprawozdawczości federalnej). Ze względu na pozycję w hierarchii, odbiorcy niższego szczebla rzadko dokonują weryfikacji takich żądań kanałem zwrotnym.
    \item \textbf{Pretexting dydaktyczny i studencki (Prof. Adrian Perrig):} Funkcja Dyrektora Studiów sprawia, że wezwania dotyczące weryfikacji zapisów na przedmioty, logowania do platform egzaminacyjnych czy potwierdzania tematów prac dyplomowych kierowane do studentów i doktorantów charakteryzują się niemal stuprocentową wiarygodnością operacyjną.
    \item \textbf{Pretexting administracyjno-proceduralny (Barbara Pfändner):} Pozycja dyrektor zarządzającej stanowi idealny punkt wyjścia do kierowania zapytań o status rozliczeń, deklaracje podatkowe, audyty sprzętowe lub weryfikację uprawnień do gabinetów i laboratoriów do pracowników administracji i sekretariatów.
    \item \textbf{Wektor wsparcia teleinformatycznego (Stefan Walter):} Zidentyfikowanie obecności szefa grupy IT (ISGINF) bezpośrednio w ścisłym zarządzie wydziału dowodzi wysokiej integracji procedur technicznych z procesami decyzyjnymi. Daje to podstawę do wysoce wiarygodnych ataków vishingowych i spear-phishingowych podszywających się pod bezpośrednie polecenia szefa IT (np. wymuszenie pilnej zmiany poświadczeń nethz lub aktualizacji kluczy SSH).
\end{enumerate}

\subsection{Instytuty badawcze i laboratoria wydziałowe}\label{subsec:instytuty-laboratoria}
% TO-DO:
% 1. Wymienić i scharakteryzować instytuty D-INFK (np. Systems, InfSec, IML, Pervasive, TCS, IVC, IPL).
% 2. Opisać centrum cyberbezpieczeństwa ZISC (współpraca z sektorem przemysłowym/bankowym).
% 3. Zwrócić uwagę na autonomię instytutów w doborze narzędzi i serwerów.

\subsection{Struktura i zadania grupy wsparcia IT (ISG D-INFK)}\label{subsec:administracja-isg}
% TO-DO:
% 1. Opisać rolę IT Service Group (ISG D-INFK) jako autonomicznego ośrodka IT wydziału.
% 2. Wskazać zadania: zarządzanie serwerownią CAB, klastrami HPC, helpdeskiem (helpdesk@inf.ethz.ch),
%    zarządzaniem stacjami roboczymi (Ubuntu LTS, Windows) i kontami nethz.
% 3. Wskazać ISG jako centralny punkt zaufania technicznego (kluczowe dla vishingu w Rozdziale 6).

\section{Otoczenie fizyczne i infrastruktura kampusowa}\label{sec:lokalizacja}

\subsection{Lokalizacja i topologia budynków wydziałowych}\label{subsec:budynki}
% TO-DO:
% 1. Opisać lokalizację budynków na kampusie ETH Zentrum (Zurych).
% 2. Skupić się na budynku CAB (Universitätstrasse 6) jako sercu wydziału (Dziekanat, ISG pokój CAB E 41.1,
%    serwerownie wydziałowe, laboratoria studenckie).
% 3. Wymienić budynki pomocnicze (np. CHN, CNB, OAT).

\subsection{Strefy dostępu i zabezpieczenia fizyczne budynków}\label{subsec:bezpieczenstwo-fizyczne}
% TO-DO:
% 1. Zestawić strefy ogólnodostępne w dzień (kantyna, korytarze, audytoria) vs strefy chronione (24/7).
% 2. Opisać system kart zbliżeniowych ETH Card (np. technologie Legic / Mifare DESFire badane w dokumentacji kampusowej).
% 3. Przeanalizować wektor fizyczny (wejście za pracownikiem, strefy dostawcze/kantyna CAB) vs procedury zamykania laboratoriów.

\section{Analiza finansowo-gospodarcza i przetargowa}\label{sec:finanse-gospodarka}

\subsection{Budżet operacyjny i struktura finansowania}\label{subsec:budzet-sprawozdania}
% TO-DO:
% 1. Zwięźle podsumować budżet ETH (>1.8 mld CHF) i mechanizm alokacji federalnej dla D-INFK.
% 2. Wskazać podział na środki podstawowe, granty zewnętrzne i kontrakty przemysłowe.

\subsection{Granty badawcze i finansowanie zewnętrzne (SNSF)}\label{subsec:granty-badawcze}
% TO-DO:
% 1. Kwerenda w bazie SNSF (data.snf.ch) dla dyscypliny Computer Science powiązana z profesorami D-INFK.
% 2. Zestawienie 2-3 realnych grantów badawczych (tematyka np. TEE/SGX, obliczenia HPC, uczenie maszynowe, sieci; kwoty w CHF).
% 3. Actionable Intelligence: powiązanie tematyki zidentyfikowanego grantu ze scenariuszem ataku w Rozdziale 6.

\subsection{Zamówienia publiczne IT i dostawcy (simap.ch)}\label{subsec:zamowienia-simap}
% TO-DO:
% 1. Wskazówka: szukać na simap.ch zamówień IT powiązanych z ETH Zurich (sprzęt np. serwery obliczeniowe, macierze, przełączniki).
% 2. Tabela 2-3 zidentyfikowanych postępowań przetargowych lub wykazanie centralizacji zakupów sprzętu przez jednostki nadrzędne ETH/BBL.
% 3. Zidentyfikować dostawców technologii jako wektor ryzyka w łańcuchu dostaw (Supply Chain).

\section{Relacje gospodarcze i ekosystem innowacji}\label{sec:partnerstwa-gospodarcze}

\subsection{Partnerstwa technologiczne i konsorcjum ZISC}\label{subsec:partnerzy-przemyslowi}
% TO-DO:
% 1. Opisać współpracę z centrami R&D w Zurychu (np. Google Zurich, Microsoft Mixed Reality Lab, IBM Research).
% 2. Wskazać konsorcjum ZISC (partnerzy bankowi i technologiczni).
% 3. Wyjaśnić wpływ obecności gości biznesowych na obniżenie czujności personelu w CAB.

\subsection{Spółki spin-off wywodzące się z D-INFK}\label{subsec:spinoffy}
% TO-DO:
% 1. Wymienić rynkowe spółki spin-off wywodzące się z D-INFK (np. Scandit, Beekeeper, DeepCode/Snyk).
% 2. Zinterpretować transfer technologii i współdzielenie narzędzi między laboratoriami a rynkiem.
